\documentclass[12pt,a4paper]{article}
\usepackage[T1]{fontenc}
\usepackage[margin=1in]{geometry}
\usepackage{graphicx}
\usepackage{amsmath}
\usepackage{amssymb}
\usepackage{listings}
\usepackage{xcolor}
\usepackage{hyperref}
\usepackage{fancyhdr}
\usepackage{float}
\usepackage{array}
\usepackage{booktabs}
\usepackage{multirow}
\usepackage{tikz}
\usepackage{circuitikz}

% Code formatting
\lstset{
    language=Verilog,
    basicstyle=\ttfamily\small,
    keywordstyle=\color{blue},
    commentstyle=\color{gray},
    stringstyle=\color{red},
    breaklines=true,
    showstringspaces=false,
    tabsize=4,
    frame=single,
    rulecolor=\color{black!30}
}

\title{\textbf{Wishbone Bus Protocol Specification}}
\author{Wishbone Open Core Bus Architecture}
\date{Revision 1.0 \\ \today}

\pagestyle{fancy}
\fancyhf{}
\rhead{Wishbone Bus Protocol}
\lhead{Specification Document}
\cfoot{\thepage}

\begin{document}

\maketitle

\newpage
\tableofcontents
\newpage

%=====================================================
\section{Introduction}
%=====================================================

\subsection{Overview}
The Wishbone Bus Protocol is a simple, yet flexible open-source bus architecture designed for System-on-Chip (SoC) integration. It provides a standardized interface for connecting master and slave peripherals within a digital system. This specification defines a detailed architecture for an 8-bit Wishbone bus implementation with synchronous operation.

\subsection{Purpose and Scope}
This document specifies the complete behavior, signal definitions, and operational protocols for a Wishbone-based system. The specification covers:
\begin{itemize}
    \item Signal definitions and timing characteristics
    \item Master and slave interface specifications
    \item Transaction protocols and handshaking mechanisms
    \item Memory access operations (read and write)
    \item Timing diagrams and signal relationships
    \item Implementation guidelines for both masters and slaves
\end{itemize}

\subsection{Design Goals}
\begin{enumerate}
    \item \textbf{Simplicity}: Minimal signal count and straightforward handshaking
    \item \textbf{Flexibility}: Support for multiple masters and slaves on a single bus
    \item \textbf{Scalability}: Extensible to larger data widths and address spaces
    \item \textbf{Synchronous Operation}: Clock-based timing for predictable behavior
    \item \textbf{Ease of Implementation}: Simple state machines for master and slave logic
\end{enumerate}

%=====================================================
\section{System Architecture}
%=====================================================

\subsection{Bus Topology}
The Wishbone bus uses a shared bus architecture where:
\begin{itemize}
    \item Multiple master devices can initiate transactions
    \item Multiple slave devices can respond to master requests
    \item A central arbiter (not specified in this document) arbitrates master access
    \item All signals are synchronized to a common clock signal
\end{itemize}

\subsection{Data Width and Address Space}
\begin{table}[H]
    \centering
    \begin{tabular}{|l|c|}
        \hline
        \textbf{Parameter} & \textbf{Value} \\
        \hline
        Data Width (DAT) & 8 bits \\
        Address Width (ADR) & 8 bits \\
        Maximum Addressable Memory & 256 locations (0x00 to 0xFF) \\
        Clock Frequency & Application Dependent \\
        \hline
    \end{tabular}
    \caption{System Parameters}
\end{table}

%=====================================================
\section{Signal Definitions}
%=====================================================

\subsection{Global Signals}

\begin{table}[H]
    \centering
    \begin{tabular}{|p{2cm}|p{1.5cm}|p{5cm}|}
        \hline
        \textbf{Signal Name} & \textbf{Width} & \textbf{Description} \\
        \hline
        CLK & 1 bit & Clock signal. All state changes occur on rising edges. \\
        RST & 1 bit & Active-high asynchronous reset. Resets all modules to initial state. \\
        \hline
    \end{tabular}
    \caption{Global Signals}
\end{table}

\subsection{Master Interface Signals}

\begin{table}[H]
    \centering
    \begin{tabular}{|p{2cm}|p{1.5cm}|p{1cm}|p{4cm}|}
        \hline
        \textbf{Signal} & \textbf{Width} & \textbf{Dir.} & \textbf{Description} \\
        \hline
        CYC & 1 bit & Out & Cycle signal. Asserted by master to indicate an active bus cycle. \\
        STB & 1 bit & Out & Strobe signal. Asserted to indicate valid address and data. \\
        WE & 1 bit & Out & Write Enable. High for write operations, low for read operations. \\
        ADR & 8 bits & Out & Address bus. Specifies the target memory address. \\
        DAT\_O & 8 bits & Out & Data Output. Write data from master to slave. \\
        ACK & 1 bit & In & Acknowledge signal. Asserted by slave to confirm transaction completion. \\
        DAT\_I & 8 bits & In & Data Input. Read data from slave to master. \\
        BUSY & 1 bit & Out & Busy flag. Indicates master is in a transaction. \\
        START & 1 bit & In & Start signal. Request from external source to initiate a transaction. \\
        ADR\_IN & 8 bits & In & Input address for the transaction. \\
        DAT\_IN & 8 bits & In & Input write data for the transaction. \\
        \hline
    \end{tabular}
    \caption{Master Interface Signals}
\end{table}

\subsection{Slave Interface Signals}

\begin{table}[H]
    \centering
    \begin{tabular}{|p{2cm}|p{1.5cm}|p{1cm}|p{4cm}|}
        \hline
        \textbf{Signal} & \textbf{Width} & \textbf{Dir.} & \textbf{Description} \\
        \hline
        CYC & 1 bit & In & Cycle signal from master. \\
        STB & 1 bit & In & Strobe signal from master. \\
        WE & 1 bit & In & Write Enable signal from master. \\
        ADR & 8 bits & In & Address bus from master. \\
        DAT\_O & 8 bits & In & Data from master (write data). \\
        ACK & 1 bit & Out & Acknowledge signal to master. \\
        DAT\_I & 8 bits & Out & Data to master (read data). \\
        \hline
    \end{tabular}
    \caption{Slave Interface Signals}
\end{table}

\subsection{Signal Descriptions}

\subsubsection{CYC (Cycle)}
\begin{itemize}
    \item \textbf{Direction}: Master Output / Slave Input
    \item \textbf{Width}: 1 bit
    \item \textbf{Timing}: Asserted when a bus cycle begins, held throughout the cycle, deasserted when the cycle ends
    \item \textbf{Purpose}: Indicates to all slaves that a valid bus cycle is in progress
    \item \textbf{Protocol}: The master asserts CYC at the start of a transaction and keeps it asserted until the slave acknowledges with ACK
\end{itemize}

\subsubsection{STB (Strobe)}
\begin{itemize}
    \item \textbf{Direction}: Master Output / Slave Input
    \item \textbf{Width}: 1 bit
    \item \textbf{Timing}: Asserted synchronously with CYC
    \item \textbf{Purpose}: Indicates that the address (ADR) and write data (DAT\_O) are valid
    \item \textbf{Protocol}: Only the slave(s) selected by the address should respond when both CYC and STB are high
\end{itemize}

\subsubsection{WE (Write Enable)}
\begin{itemize}
    \item \textbf{Direction}: Master Output / Slave Input
    \item \textbf{Width}: 1 bit
    \item \textbf{Timing}: Valid during CYC and STB
    \item \textbf{Purpose}: Specifies the transaction type
    \item \textbf{Values}:
    \begin{itemize}
        \item WE = 1: Write operation (master to slave data transfer)
        \item WE = 0: Read operation (slave to master data transfer)
    \end{itemize}
\end{itemize}

\subsubsection{ADR (Address Bus)}
\begin{itemize}
    \item \textbf{Direction}: Master Output / Slave Input
    \item \textbf{Width}: 8 bits
    \item \textbf{Timing}: Valid when CYC and STB are high
    \item \textbf{Purpose}: Specifies the target memory location or register address
    \item \textbf{Range}: 0x00 to 0xFF (256 addressable locations)
\end{itemize}

\subsubsection{DAT\_O (Data Output from Master)}
\begin{itemize}
    \item \textbf{Direction}: Master Output / Slave Input
    \item \textbf{Width}: 8 bits
    \item \textbf{Timing}: Valid when CYC, STB, and WE are high
    \item \textbf{Purpose}: Carries write data from master to slave
    \item \textbf{Hold Time}: Data is held stable until ACK is asserted
\end{itemize}

\subsubsection{DAT\_I (Data Input to Master)}
\begin{itemize}
    \item \textbf{Direction}: Slave Output / Master Input
    \item \textbf{Width}: 8 bits
    \item \textbf{Timing}: Valid when CYC and STB are high, and WE is low (read operation)
    \item \textbf{Purpose}: Carries read data from slave to master
    \item \textbf{Behavior}: Slave must drive this bus with the requested data at the target address
\end{itemize}

\subsubsection{ACK (Acknowledge)}
\begin{itemize}
    \item \textbf{Direction}: Slave Output / Master Input
    \item \textbf{Width}: 1 bit
    \item \textbf{Timing}: Asserted by slave in response to a valid transaction
    \item \textbf{Purpose}: Signals transaction completion to the master
    \item \textbf{Protocol}: 
    \begin{itemize}
        \item Asserted when CYC and STB conditions are met
        \item Held high for one clock cycle
        \item Signals both read and write completion
    \end{itemize}
\end{itemize}

%=====================================================
\section{Transaction Protocol}
%=====================================================

\subsection{Transaction Types}

\subsubsection{Write Transaction}
A write transaction transfers data from the master to the slave at a specified address.

\textbf{Sequence}:
\begin{enumerate}
    \item Master asserts START signal with ADR\_IN (target address) and DAT\_IN (write data)
    \item On next clock edge, master sets:
    \begin{itemize}
        \item CYC = 1
        \item STB = 1
        \item WE = 1 (write mode)
        \item ADR = ADR\_IN
        \item DAT\_O = DAT\_IN
        \item BUSY = 1
    \end{itemize}
    \item Slave observes CYC=1, STB=1, WE=1 and valid ADR
    \item Slave captures DAT\_O and writes to memory at address ADR
    \item Slave asserts ACK = 1
    \item Master observes ACK = 1 on next clock edge
    \item Master deasserts:
    \begin{itemize}
        \item CYC = 0
        \item STB = 0
        \item BUSY = 0
    \end{itemize}
    \item Slave deasserts ACK = 0
    \item Transaction completes
\end{enumerate}

\subsubsection{Read Transaction}
A read transaction transfers data from the slave to the master from a specified address.

\textbf{Sequence}:
\begin{enumerate}
    \item Master asserts START signal with ADR\_IN (target address), DAT\_IN = 0
    \item On next clock edge, master sets:
    \begin{itemize}
        \item CYC = 1
        \item STB = 1
        \item WE = 0 (read mode)
        \item ADR = ADR\_IN
        \item BUSY = 1
    \end{itemize}
    \item Slave observes CYC=1, STB=1, WE=0 and valid ADR
    \item Slave reads memory at address ADR and drives DAT\_I with the value
    \item Slave asserts ACK = 1
    \item Master observes ACK = 1 and captures DAT\_I on the same clock edge
    \item Master deasserts:
    \begin{itemize}
        \item CYC = 0
        \item STB = 0
        \item BUSY = 0
    \end{itemize}
    \item Slave deasserts ACK = 0
    \item Transaction completes
\end{enumerate}

\subsection{Transaction Timing}

\textbf{Key Timing Characteristics}:
\begin{itemize}
    \item \textbf{Single-Cycle Transactions}: Both read and write complete in exactly 2 clock cycles from START assertion
    \item \textbf{Setup Time}: Address and data must be stable before clock edge
    \item \textbf{Hold Time}: Address and data are held through ACK assertion
    \item \textbf{Clock-to-ACK Delay}: Slave ACK appears synchronously on clock edge following valid strobe signals
\end{itemize}

\subsection{Protocol Rules}

\subsubsection{For Masters}
\begin{enumerate}
    \item Master must maintain CYC and STB asserted until ACK is received
    \item Address (ADR) and write data (DAT\_O) must be stable while CYC and STB are asserted
    \item Master must deassert CYC and STB within one clock cycle of observing ACK
    \item BUSY signal must accurately reflect transaction state
    \item START signal should be asserted for one clock cycle only
\end{enumerate}

\subsubsection{For Slaves}
\begin{enumerate}
    \item Slave must respond with ACK when both CYC and STB are high and addressing matches
    \item For write transactions, slave must capture DAT\_O and store at address ADR
    \item For read transactions, slave must immediately drive DAT\_I with data from address ADR
    \item Slave must deassert ACK after one clock cycle
    \item Slave must not respond to transactions when CYC or STB are low
\end{enumerate}

%=====================================================
\section{Timing Diagrams}
%=====================================================

\subsection{Write Transaction Timing}

\begin{figure}[H]
\centering
\begin{verbatim}
Clock      ___     ___     ___     ___     ___     ___     ___
           ___|   |___|   |___|   |___|   |___|   |___|   |___|
START      ___
           ____|_______________________
           
CYC        ___
           ____|_________________|___________
           
STB        ___
           ____|_________________|___________
           
WE         ___
           ____________________________________
           
ADR        ___     [Valid Address]     [Valid Address]
           ____|_________________|___________
           
DAT_O      ___     [Write Data]        [Write Data]
           ____|_________________|___________
           
ACK        ___
           ______________________|___________
           
BUSY       ___
           ____|_________________|___________

           Cycle 0  1      2      3      4
           
\end{verbatim}
\caption{Write Transaction Timing Diagram}
\end{figure}

\subsection{Read Transaction Timing}

\begin{figure}[H]
\centering
\begin{verbatim}
Clock      ___     ___     ___     ___     ___     ___     ___
           ___|   |___|   |___|   |___|   |___|   |___|   |___|
START      ___
           ____|_______________________
           
CYC        ___
           ____|_________________|___________
           
STB        ___
           ____|_________________|___________
           
WE         ___
           ____________________________________
           
ADR        ___     [Valid Address]     [Valid Address]
           ____|_________________|___________
           
DAT_I      ___     [Read Data]         [Read Data]
           ____|_________________|___________
           
ACK        ___
           ______________________|___________
           
BUSY       ___
           ____|_________________|___________

           Cycle 0  1      2      3      4
           
\end{verbatim}
\caption{Read Transaction Timing Diagram}
\end{figure}

%=====================================================
\section{Implementation Specifications}
%=====================================================

\subsection{Master Implementation}

\subsubsection{State Machine}
The master employs a simple two-state state machine:

\textbf{IDLE State}:
\begin{itemize}
    \item CYC = 0, STB = 0, BUSY = 0
    \item Waits for START signal assertion
    \item On START: transition to BUSY state
\end{itemize}

\textbf{BUSY State}:
\begin{itemize}
    \item CYC = 1, STB = 1, BUSY = 1
    \item Holds transaction signals until ACK received
    \item On ACK: deassert signals and return to IDLE state
\end{itemize}

\subsubsection{Register Operations}
\begin{lstlisting}[caption=Master State Machine Logic]
always @(posedge clk or posedge rst)
begin
    if (rst)
    begin
        cyc <= 0;
        stb <= 0;
        we <= 0;
        adr <= 0;
        dat_o <= 0;
        busy <= 0;
    end
    else
    begin
        // Transition from IDLE to BUSY
        if (!busy && start)
        begin
            cyc <= 1;
            stb <= 1;
            we <= 1;
            adr <= adr_in;
            dat_o <= dat_in;
            busy <= 1;
        end
        // Transition from BUSY to IDLE
        else if (busy && ack)
        begin
            cyc <= 0;
            stb <= 0;
            busy <= 0;
        end
    end
end
\end{lstlisting}

\subsubsection{Operation Summary}
\begin{itemize}
    \item All state changes occur on clock rising edges
    \item Reset is asynchronous and active-high
    \item Master initiates transactions via START signal
    \item BUSY flag prevents overlapping transactions
    \item Data width is 8 bits for simplicity
\end{itemize}

\subsection{Slave Implementation}

\subsubsection{Memory Architecture}
The slave contains an internal 256-byte addressable memory:
\begin{lstlisting}[caption=Slave Memory Declaration]
reg [7:0] memory [0:255];
\end{lstlisting}

\subsubsection{State Machine}
The slave employs a request-response protocol:

\textbf{IDLE State}:
\begin{itemize}
    \item ACK = 0
    \item Monitors CYC and STB signals
    \item On CYC=1 and STB=1: process transaction
\end{itemize}

\textbf{RESPOND State}:
\begin{itemize}
    \item ACK = 1 for one clock cycle
    \item For write: captures DAT\_O into memory[ADR]
    \item For read: drives DAT\_I with memory[ADR]
\end{itemize}

\subsubsection{Register Operations}
\begin{lstlisting}[caption=Slave State Machine Logic]
always @(posedge clk or posedge rst)
begin
    if (rst)
    begin
        ack <= 0;
    end
    else
    begin
        if (cyc && stb)
        begin
            ack <= 1;
            
            if (we)
                memory[adr] <= dat_o;   // Write operation
            else
                dat_i <= memory[adr];   // Read operation
        end
        else
            ack <= 0;
    end
end
\end{lstlisting}

\subsubsection{Operation Summary}
\begin{itemize}
    \item All transactions complete in one clock cycle
    \item Write operations are synchronous (data captured on clock edge)
    \item Read operations provide data combinatorially through DAT\_I
    \item Slave supports full 8-bit address range (0x00 to 0xFF)
    \item Slave supports all 256 memory locations
\end{itemize}

%=====================================================
\section{System Operation}
%=====================================================

\subsection{Reset Sequence}

Upon assertion of RST:
\begin{enumerate}
    \item All flip-flops in master are reset to 0
    \item All flip-flops in slave are reset to 0
    \item Bus signals (CYC, STB, WE, ACK) are deasserted
    \item Both master and slave return to IDLE state
    \item Memory contents in slave are preserved (not reset)
\end{enumerate}

\subsection{Normal Operation Flow}

\begin{enumerate}
    \item System starts with RST = 1 for at least one clock cycle
    \item RST is deasserted; both master and slave enter normal operation
    \item External stimulus provides START, ADR\_IN, and DAT\_IN to master
    \item Master begins transaction as per protocol
    \item Slave responds with data or acknowledgment
    \item Transaction completes with CYC and STB deasserted
    \item System ready for next transaction
\end{enumerate}

\subsection{Error Handling}

\subsubsection{Invalid Addresses}
In this implementation, all 256 addresses (0x00-0xFF) are valid and accessible. Out-of-range addresses wrap around due to 8-bit address bus.

\subsubsection{Bus Timeout}
This specification does not include bus timeout mechanisms. Masters must ensure:
\begin{itemize}
    \item Slave responds with ACK within finite time
    \item No deadlock conditions exist
    \item All addressed locations are valid
\end{itemize}

\subsubsection{Multiple Masters}
If multiple masters are present in the system:
\begin{itemize}
    \item An external arbiter must control which master drives the bus
    \item Only one master can assert CYC at a time
    \item Arbitration logic is outside this specification
\end{itemize}

%=====================================================
\section{Advanced Features}
%=====================================================

\subsection{Extensibility}

This specification can be extended to support:

\begin{itemize}
    \item \textbf{Wider Data Paths}: Increase DAT\_O and DAT\_I widths to 16, 32, or 64 bits
    \item \textbf{Larger Address Space}: Extend ADR width beyond 8 bits for larger memory
    \item \textbf{Burst Transactions}: Extend protocol to support multiple data transfers per cycle initiation
    \item \textbf{Byte Enables}: Add SEL (select) signals for sub-word access
    \item \textbf{Error Signaling}: Add ERR signal for fault conditions
    \item \textbf{Ready Signals}: Add RDY signal for pipelined transactions
\end{itemize}

\subsection{Performance Optimization}

\subsubsection{Burst Mode}
Extend the protocol to support multiple consecutive accesses without releasing the bus:
\begin{itemize}
    \item Keep CYC asserted across multiple transactions
    \item Perform address and data changes between transactions
    \item Only release bus when all data transfers complete
\end{itemize}

\subsubsection{Pipelined Transfers}
Allow new transactions to begin before previous ones complete:
\begin{itemize}
    \item Requires additional buffering
    \item Improves bus utilization
    \item Requires more complex arbiter logic
\end{itemize}

%=====================================================
\section{Verification and Testing}
%=====================================================

\subsection{Test Cases}

\subsubsection{Write Verification}
\begin{enumerate}
    \item Write 0xA5 to address 0x10
    \item Verify memory location 0x10 contains 0xA5
    \item Repeat for addresses 0x00, 0x80, 0xFF
\end{enumerate}

\subsubsection{Read Verification}
\begin{enumerate}
    \item Pre-load memory with known patterns
    \item Read from various addresses
    \item Verify returned data matches memory contents
    \item Test all addresses in range
\end{enumerate}

\subsubsection{Sequential Operations}
\begin{enumerate}
    \item Perform multiple read-write sequences
    \item Verify BUSY signal prevents overlapping operations
    \item Verify no data corruption
\end{enumerate}

\subsubsection{Edge Cases}
\begin{enumerate}
    \item Test address wraparound (0xFF to 0x00)
    \item Test boundary addresses (0x00, 0xFF)
    \item Test all possible 8-bit data values
    \item Test reset during operation
\end{enumerate}

\subsection{Functional Coverage}

All operations must achieve functional coverage for:
\begin{itemize}
    \item \textbf{Read Operations}: All 256 addresses read with valid data returned
    \item \textbf{Write Operations}: All 256 addresses written with data correctly stored
    \item \textbf{Reset Behavior}: Signals correctly initialized on RST
    \item \textbf{State Transitions}: All master state transitions exercised
    \item \textbf{ACK Timing}: ACK timing verified across multiple transactions
\end{itemize}

%=====================================================
\section{Pin Configuration and Connectivity}
%=====================================================

\subsection{Master Port Connections}

\begin{table}[H]
    \centering
    \begin{tabular}{|l|l|l|}
        \hline
        \textbf{Signal} & \textbf{Master} & \textbf{Slave} \\
        \hline
        CLK & Input & Input \\
        RST & Input & Input \\
        CYC & Output & Input \\
        STB & Output & Input \\
        WE & Output & Input \\
        ADR[7:0] & Output & Input \\
        DAT\_O[7:0] & Output & Input \\
        DAT\_I[7:0] & Input & Output \\
        ACK & Input & Output \\
        \hline
    \end{tabular}
    \caption{Master-Slave Signal Connectivity}
\end{table}

\subsection{Interface Matching}
\begin{itemize}
    \item All outputs from one module connect to inputs of the other
    \item Bus signals (CYC, STB, WE, ADR, DAT\_O) are point-to-point from master to slave
    \item Response signals (ACK, DAT\_I) are point-to-point from slave to master
    \item Clock and reset are distributed to all modules
\end{itemize}

%=====================================================
\section{Application Examples}
%=====================================================

\subsection{Simple Memory System}
A basic system with a single master and single slave:
\begin{itemize}
    \item Master: Microcontroller or state machine
    \item Slave: 256-byte RAM
    \item Application: Memory-mapped I/O and data storage
\end{itemize}

\subsection{Multi-Slave System}
Extension to multiple slave devices:
\begin{itemize}
    \item Single master drives up to 256 addressable slaves
    \item Each slave responds to its assigned address range
    \item Slave selection based on upper address bits
\end{itemize}

\subsection{Multi-Master System}
Multiple masters sharing a single bus:
\begin{itemize}
    \item Central arbiter grants bus access to masters
    \item Prevents address conflicts through arbitration
    \item Requires priority or round-robin arbitration
\end{itemize}

%=====================================================
\section{Design Guidelines}
%=====================================================

\subsection{Clock Domain Recommendations}

\begin{itemize}
    \item \textbf{Synchronous Design}: All Wishbone signals should operate in the same clock domain
    \item \textbf{Clock Frequency}: Determined by slowest component in the system
    \item \textbf{Reset Synchronization}: RST should be synchronized to CLK before distribution
    \item \textbf{Clock Distribution}: Use a dedicated clock tree to minimize skew
\end{itemize}

\subsection{Signal Integrity}

\begin{itemize}
    \item \textbf{Line Termination}: Implement appropriate termination for high-speed designs
    \item \textbf{Ground Reference}: Maintain clean ground planes for signal integrity
    \item \textbf{Impedance Control}: Control trace impedance to minimize reflections
    \item \textbf{Crosstalk}: Route address and data buses away from high-speed signals
\end{itemize}

\subsection{Power Considerations}

\begin{itemize}
    \item \textbf{Power Supplies}: Separate supply rails for analog and digital circuits when necessary
    \item \textbf{Decoupling}: Place bypass capacitors near all power pins
    \item \textbf{Ground Distribution}: Use multiple vias for low impedance ground connections
\end{itemize}

%=====================================================
\section{Electrical Specifications}
%=====================================================

\subsection{DC Characteristics}

\begin{table}[H]
    \centering
    \begin{tabular}{|l|c|c|}
        \hline
        \textbf{Parameter} & \textbf{Min} & \textbf{Max} \\
        \hline
        Supply Voltage (VDD) & 3.0 V & 3.6 V \\
        Logic High (VIH) & 2.0 V & VDD \\
        Logic Low (VIL) & 0 V & 0.8 V \\
        Output High (VOH) & VDD-0.2 V & VDD \\
        Output Low (VOL) & 0 V & 0.2 V \\
        \hline
    \end{tabular}
    \caption{DC Electrical Specifications (Typical for 3.3V CMOS)}
\end{table}

\subsection{AC Characteristics}

\begin{table}[H]
    \centering
    \begin{tabular}{|l|c|}
        \hline
        \textbf{Parameter} & \textbf{Value} \\
        \hline
        Setup Time (tSU) & 2 ns (typical) \\
        Hold Time (tH) & 1 ns (typical) \\
        Propagation Delay (tPD) & 3 ns (typical) \\
        Clock-to-Output (tCO) & 5 ns (typical) \\
        \hline
    \end{tabular}
    \caption{AC Timing Characteristics}
\end{table}

%=====================================================
\section{Revision History}
%=====================================================

\begin{table}[H]
    \centering
    \begin{tabular}{|c|c|l|}
        \hline
        \textbf{Revision} & \textbf{Date} & \textbf{Description} \\
        \hline
        1.0 & 2026-04-22 & Initial specification document \\
        \hline
    \end{tabular}
    \caption{Document Revision History}
\end{table}

%=====================================================
\section{References and Acknowledgments}
%=====================================================

\subsection{External Resources}
\begin{itemize}
    \item Wishbone Architecture Documentation
    \item AMBA Bus Protocol Specification
    \item IEEE Standard 1735 (Intellectual Property Protection)
\end{itemize}

\subsection{Contributors}
\begin{itemize}
    \item Wishbone Open Core Community
    \item Protocol Implementation Team
\end{itemize}

\newpage
\appendix

%=====================================================
\section{Appendix A: Verilog Module Definitions}
%=====================================================

\subsection{Master Module}

\begin{lstlisting}[caption=Complete Wishbone Master Module]
module wishbone_master(
    input clk,
    input rst,
    input ack,
    input [7:0] dat_i,
    input start,
    input [7:0] adr_in,
    input [7:0] dat_in,
    
    output reg cyc,
    output reg stb,
    output reg we,
    output reg [7:0] dat_o,
    output reg [7:0] adr,
    output reg busy
);

always @(posedge clk or posedge rst)
begin
    if (rst)
    begin
        cyc <= 0;
        stb <= 0;
        we <= 0;
        adr <= 0;
        dat_o <= 0;
        busy <= 0;
    end
    else
    begin
        if (!busy && start)
        begin
            cyc <= 1;
            stb <= 1;
            we <= 1;
            adr <= adr_in;
            dat_o <= dat_in;
            busy <= 1;
        end
        else if (busy && ack)
        begin
            cyc <= 0;
            stb <= 0;
            busy <= 0;
        end
    end
end

endmodule
\end{lstlisting}

\subsection{Slave Module}

\begin{lstlisting}[caption=Complete Wishbone Slave Module]
module wishbone_slave(
    input clk,
    input rst,
    input cyc,
    input stb,
    input we,
    input [7:0] adr,
    input [7:0] dat_o,
    
    output reg ack,
    output reg [7:0] dat_i
);

reg [7:0] memory [0:255];

always @(posedge clk or posedge rst)
begin
    if (rst)
    begin
        ack <= 0;
    end
    else
    begin
        if (cyc && stb)
        begin
            ack <= 1;
            
            if (we)
                memory[adr] <= dat_o;
            else
                dat_i <= memory[adr];
        end
        else
            ack <= 0;
    end
end

endmodule
\end{lstlisting}

%=====================================================
\section{Appendix B: Test Bench Structure}
%=====================================================

\subsection{Recommended Test Bench Organization}

\begin{lstlisting}[caption=Testbench Framework]
module wishbone_tb;

reg clk, rst;
reg [7:0] adr_in, dat_in;
reg start;

wire cyc, stb, we, ack, busy;
wire [7:0] adr, dat_o, dat_i;

wishbone_master master(
    .clk(clk), .rst(rst), .start(start),
    .adr_in(adr_in), .dat_in(dat_in),
    .cyc(cyc), .stb(stb), .we(we),
    .adr(adr), .dat_o(dat_o), .dat_i(dat_i),
    .ack(ack), .busy(busy)
);

wishbone_slave slave(
    .clk(clk), .rst(rst),
    .cyc(cyc), .stb(stb), .we(we),
    .adr(adr), .dat_o(dat_o),
    .ack(ack), .dat_i(dat_i)
);

initial begin
    // Test procedures
end

endmodule
\end{lstlisting}

%=====================================================
\section{Appendix C: Address Decoding Examples}
%=====================================================

\subsection{Multi-Slave Address Decoding}

For systems with multiple slaves, use address decoding:

\begin{lstlisting}[caption=Address Decoder Example]
// Slave 0: addresses 0x00-0x3F
// Slave 1: addresses 0x40-0x7F
// Slave 2: addresses 0x80-0xBF
// Slave 3: addresses 0xC0-0xFF

wire slave0_sel = (adr[7:6] == 2'b00);
wire slave1_sel = (adr[7:6] == 2'b01);
wire slave2_sel = (adr[7:6] == 2'b10);
wire slave3_sel = (adr[7:6] == 2'b11);
\end{lstlisting}

%=====================================================
\section{Appendix D: Glossary}
%=====================================================

\begin{table}[H]
    \centering
    \begin{tabular}{|p{3cm}|p{7cm}|}
        \hline
        \textbf{Term} & \textbf{Definition} \\
        \hline
        ACK & Acknowledge signal from slave indicating transaction completion \\
        ADR & Address bus carrying memory or register location \\
        Bus Cycle & Complete transaction from CYC assertion to ACK \\
        CYC & Cycle signal indicating active bus transaction \\
        DAT & Data bus (DAT\_O for output, DAT\_I for input) \\
        Master & Initiator of bus transactions \\
        Slave & Responder to bus transactions \\
        STB & Strobe signal indicating valid address and data \\
        WE & Write Enable signal (1=write, 0=read) \\
        \hline
    \end{tabular}
    \caption{Glossary of Terms}
\end{table}

%=====================================================
\end{document}
